\documentclass[11pt]{article}
\usepackage[margin=1in]{geometry}
\usepackage[shortlabels]{enumitem}
\usepackage{amsmath, amssymb, amsthm}
\usepackage{newpxtext, newpxmath}

\newtheorem{thm}{Theorem}[subsection]
\newtheorem{lemma}[thm]{Lemma}
\newtheorem{corollary}[thm]{Corollary}
\newtheorem{prop}[thm]{Proposition}

\theoremstyle{definition}
\newtheorem{definition}[thm]{Definition}
\newtheorem{example}[thm]{Example}

\theoremstyle{remark}
\newtheorem{remark}[thm]{Remark}

\title{\textbf{Discrete Mathematics}}
\author{Alex Davison}

\begin{document}
\maketitle

\tableofcontents

\section{Introduction}

These notes on discrete mathematics cover material that intersects both foundational mathematics and computer science. Logic and proof are paramount, so results will be justified with proof. Intuition and insight is highly prioritized over memorization and mechanical procedures (``plug and chug'').

Exactly what counts as ``discrete mathematics'' is a bit of a loose matter. Calculus and geometry seem ``continuous'' and so are not topics here; functions, sets, logic, integers, counting problems, graphs are topics of interest instead.

\section{Logic and sets}

Mathematics is grounded in \emph{formal, deductive logic}. Results in mathematics are derived from logical proof. The typical treatment of mathematics in grade school only ambiently assumes this logic and rarely discusses any proofs, but mathematics for the pure mathematician is almost all about proof. This section discusses logic and proof at a high level, and introduces the fundamental class of mathematical objects, sets.

\subsection{Logic and proof}

Logic operates over \emph{propositions}, which are mathematical statements that hold some truth value semantically, either ``true'' or ``false.'' Euclid's result ``there exist infinitely many primes'' is an example of a true proposition; the claim that ``there exists a rational number $x$ such that $x^2 = 2$'' is an example of a false proposition. The truthhood of a proposition is unknown until a \emph{proof} is supplied indicating whether it is true or false.

Suppose $A$ and $B$ are propositions. There are a few \emph{logical connectives} allowing you to make compound propositions:
\begin{itemize}
    \item $A \wedge B$ (conjunction)
    \item $A \vee B$ (disjunction)
    \item $A \implies B$ (implication)
    \item $A \iff B$ (equivalence)
\end{itemize}

Conjunction $A \wedge B$ means ``$A$ and $B$,'' which is true precisely when both $A$ and $B$ are true. Disjunction $A \vee B$ means ``$A$ or $B$,'' which is true if either $A$ or $B$ is true, and false only if both are false. To prove $A \wedge B$, you obviously must prove each of $A$ and $B$. To prove $A \vee B$, it suffices to prove one of $A$ or $B$.

Implication $A \implies B$ can be read a few different ways:
\begin{itemize}
    \item $A$ implies $B$,
    \item $A$ only if $B$,
    \item $A$ is sufficient for $B$,
    \item $B$ is implied by $A$,
    \item $B$ if $A$,
    \item $B$ is necessary for $A$.
\end{itemize}
Perhaps the best way to think of a proposition $A \implies B$ is as a guarantee: Whenever $A$ is true, it is guaranteed that $B$ follows. But a guarantee need not promise anything in the event that $A$ is false. So the statement ``if unicorns are real, then pigs can fly'' is perfectly true logically, since the hypothesis that unicorns are real is false in the first place. Propositions like these are called \emph{vacuously true}; the fact that any proposition is implied by a falsity is called the \emph{principle of explosion} or \emph{ex falso quodlibet}.

If we know that $A$ is true and we know $A \implies B$, then we may conclude $B$. This obvious fact is called \emph{modus ponens}, the foundational inference rule in logic. To prove $A \implies B$, you may temporarily assume $A$ to be true (``Suppose that $A$...''), and then derive that $B$ is true.

If you assume $A$ is true and arrive at a \emph{contradiction}, denoted by [TODO]

\subsection{Sets and set notation}

\begin{definition}
    A \emph{set} is a collection of objects sharing some common property. If $\phi(x)$ denotes some property over $x$, then we denote by $\{x \mid \phi(x)\}$ the set of all $x$ such that $\phi(x)$ holds. The vertical bar is read ``such that.'' If an object $a$ is in a set $A$, we say that $a$ is a \emph{member} or \emph{element} of $A$ and write $a \in A$ (``$x$ is in $A$''). Thus, $a \in \{x \mid \phi(x)\}$ means that $\phi(a)$ holds.
\end{definition}

\begin{example}
    Let's say we want to talk about the set of even integers. Suppose we already defined the set of integers, typically denoted by $\mathbb{Z}$, and we already defined the property of evenness. Then we can write $\{x \in \mathbb{Z} \mid x \text{ is even}\}$, read ``the set of integers $x \in \mathbb{Z}$ such that $x$ is even.''
\end{example}

\begin{definition}
    Suppose $A$ and $B$ are sets. Some set operations:
    \begin{itemize}
        \item $A \cup B$ denotes the \emph{union} of $A$ and $B$, i.e., the set of both $A$'s elements and $B$'s elements.
        \item $A \cap B$ denotes the \emph{intersection} of $A$ and $B$, i.e., the set of elements shared across both $A$ and $B$.
        \item $A \smallsetminus B$ denoted the \emph{set difference} $A$ minus $B$, i.e., the set $A$ but with elements of $B$ removed.
    \end{itemize}
\end{definition}

\begin{prop}[De Morgan's laws]
    $A \smallsetminus (B \cup C) = (A \smallsetminus B) \cap (A \smallsetminus C)$ and $A \smallsetminus (B \cap C) = (A \smallsetminus B) \cup (A \smallsetminus C)$.
\end{prop}

\end{document}